\section{Introduction}
\label{sec:introduction}

Language models may retain knowledge that their developers later wish to
remove, for example because it is private, copyrighted, obsolete, or unsafe.
Post-hoc concept erasure attempts to remove such knowledge from model
parameters while preserving unrelated capabilities. A successful method should
do more than change answers to a familiar prompt: its effect should persist
across task formats and should not merely hide accessible knowledge.

Most evaluations compare an erased model only with its original version. This
establishes that an intervention changed measured behavior, but not that the
result resembles genuine absence of learning. An edit can be too weak, too
strong, or damaging in a qualitatively different way while still obtaining a
large efficacy score. Retraining on data from which the target concept has been
removed supplies a more natural reference, but the expense of training and the
difficulty of identifying concept-related data have usually made this baseline
unavailable.

We exploit LMEnt's entity-annotated corpus and open training pipeline
\citep{gottesman2025lment} to construct matched OLMo-2 1B models
\citep{walsh2025olmo2}. A shared control is trained on the complete corpus,
while two matched models receive no loss on chunks identified with Ancient Rome
or Baseball. Because entity filtering cannot guarantee removal of every
indirect mention, we call these models \emph{concept-excluded}; they are our
operational never-trained references.

Our primary method is EMBER \citep{suslik2026ember}, which removes
concept-associated sparse features from selected token embeddings. RMU
\citep{li2024wmdp} and semi-nonnegative matrix factorization (SNMF)
\citep{shafran2026snmf} provide MLP-editing comparisons. We ask:

\begin{quote}
\small
How closely do post-hoc erasure methods reproduce a matched model trained
without learning signal from concept-associated text?
\end{quote}

No single distance answers this question. We therefore compare the models on
concept questions, held-out concept text, unrelated controls, and parameter
displacements. The present EMBER results disagree across both concepts and
instruments. On Ancient Rome, EMBER overshoots concept exclusion and damages a
tail of passages that the excluded model degrades only mildly. On Baseball, it
reaches only part of the exclusion effect on held-out text; its clear
question-likelihood reduction cannot be expressed as an overshoot because the
excluded-model question effect is within the between-twin floor. The completed
Baseball parameter comparison
shows almost no alignment between the embedding edit and the displacement
caused by excluding data during training.

Our contributions are threefold:

\begin{itemize}
  \item We construct controlled concept-excluded references for two concepts
  while holding initialization, batch order, optimizer, and training duration
  fixed.
  \item We evaluate erasure as similarity to those references across
  behavioral, distributional, parameter-space, and specificity dimensions.
  \item We show that these dimensions can disagree, and that a saturated
  accuracy-based selection metric can conceal collateral effects and favor an
  unnecessarily strong edit.
\end{itemize}

\placeholder{Revise the result preview and contribution list after RMU and SNMF
are complete; keep EMBER as the main analysis.}
